\chapter{Primitive Roots}
\intro{
    Euler's theorem says that $a^{\phi(m)} \equiv 1 \pmod m$ for $(a,m) = 1$. The smallest exponent with this property is the order of $a$, and a primitive root is an element whose order is as large as possible, namely $\phi(m)$.

    We prove that a primitive root exists modulo $m$ exactly when $m = 1, 2, 4, p^k$ or $2p^k$ with $p$ an odd prime. In other words, these are the moduli for which $(\mathbb{Z}/m\mathbb{Z})^\times$ is cyclic. The order $\ord_q(p)$ appears in the analytic part when we factor $\prod_\chi L(s,\chi)$.
}
\section{Order of an Element}

\begin{definition}[Order]\label{def:order}
    Let $m \ge 1$ and $(a,m) = 1$. The order of $a$ modulo $m$, denoted $\ord_m(a)$, is the smallest integer $k \ge 1$ such that $a^k \equiv 1 \pmod m$. It exists by \cref{thm:euler_theorem}.
\end{definition}

\begin{proposition}\label{prop:order_divides}
    Let $(a,m) = 1$ and $h = \ord_m(a)$. For $n \ge 0$, we have $a^n \equiv 1 \pmod m$ if and only if $h \mid n$. In particular, $\ord_m(a) \mid \phi(m)$, and $a^i \equiv a^j \pmod m$ if and only if $i \equiv j \pmod h$.
\end{proposition}
\begin{proof}
    Write $n = hq + r$ with $0 \le r < h$. Then $a^n = (a^h)^q a^r \equiv a^r \pmod m$, and by minimality of $h$, $a^r \equiv 1$ if and only if $r = 0$. The second claim follows from \cref{thm:euler_theorem}. For the last one, if $i \ge j$, then $a^i \equiv a^j$ is equivalent to $a^{i-j} \equiv 1$ by \cref{cor:cancellation}.
\end{proof}

\begin{proposition}\label{prop:order_of_power}
    Let $(a,m) = 1$, $h = \ord_m(a)$ and $k \ge 1$. Then
    \begin{equation*}
        \ord_m(a^k) = \frac{h}{(h,k)}.
    \end{equation*}
\end{proposition}
\begin{proof}
    By \cref{prop:order_divides}, $(a^k)^j \equiv 1$ if and only if $h \mid kj$. Dividing by $(h,k)$ and using \cref{thm:gauss_lemma}, this happens if and only if $\frac{h}{(h,k)} \mid j$. The smallest such $j \ge 1$ is $h/(h,k)$.
\end{proof}

\begin{lemma}\label{lem:order_product}
    Let $(ab,m) = 1$. If $\ord_m(a) = h$, $\ord_m(b) = k$ and $(h,k) = 1$, then $\ord_m(ab) = hk$.
\end{lemma}
\begin{proof}
    Clearly $(ab)^{hk} \equiv 1$. Conversely, suppose $(ab)^n \equiv 1$. Raising to the power $h$ gives $b^{hn} \equiv 1$, so $k \mid hn$ and therefore $k \mid n$. Symmetrically $h \mid n$, and since $(h,k) = 1$ we get $hk \mid n$.
\end{proof}

\begin{definition}[Primitive root]\label{def:primitive_root}
    An integer $g$ with $(g,m) = 1$ is a primitive root modulo $m$ if $\ord_m(g) = \phi(m)$.
\end{definition}
By \cref{prop:order_divides}, if $g$ is a primitive root, the powers $g^0, g^1, \dots, g^{\phi(m)-1}$ are pairwise incongruent, so they form a reduced residue system modulo $m$. Thus $g$ is a primitive root exactly when it generates the group $(\mathbb{Z}/m\mathbb{Z})^\times$, which is then cyclic. By \cref{prop:order_of_power}, $g^k$ is again a primitive root if and only if $(k, \phi(m)) = 1$, so when primitive roots exist there are exactly $\phi(\phi(m))$ of them modulo $m$.

\section{Primitive Roots Modulo a Prime}

\begin{theorem}[Lagrange]\label{thm:lagrange_roots}
    Let $p$ be prime and $f \in \mathbb{Z}[x]$ a polynomial of degree $n \ge 0$ whose leading coefficient is not divisible by $p$. Then the congruence $f(x) \equiv 0 \pmod p$ has at most $n$ solutions modulo $p$.
\end{theorem}
\begin{proof}
    We argue by induction on $n$. For $n = 0$, $f$ is a constant not divisible by $p$ and there is no solution. Let $n \ge 1$ and suppose $f(a) \equiv 0 \pmod p$. Dividing $f$ by $x - a$ gives $f(x) = (x - a)g(x) + f(a)$, where $g$ has degree $n - 1$ and the same leading coefficient as $f$. If $f(b) \equiv 0$ with $b \not\equiv a$, then $p \mid (b-a)g(b)$, so $g(b) \equiv 0$ by \cref{cor:euclid_lemma}. By induction there are at most $n - 1$ such $b$, hence at most $n$ roots of $f$ in total.
\end{proof}
The hypothesis that $p$ is prime is essential: $x^2 - 1$ has four roots modulo $8$.

\begin{lemma}\label{lem:roots_xd_minus_1}
    If $d \mid p - 1$, the congruence $x^d \equiv 1 \pmod p$ has exactly $d$ solutions modulo $p$.
\end{lemma}
\begin{proof}
    Write $p - 1 = de$. Then
    \begin{equation*}
        x^{p-1} - 1 = (x^d - 1)h(x), \qquad h(x) = x^{d(e-1)} + x^{d(e-2)} + \dots + x^d + 1.
    \end{equation*}
    By \cref{cor:fermat}, $x^{p-1} - 1$ has the $p - 1$ roots $1, \dots, p-1$. Each of them is a root of $x^d - 1$ or of $h$ by \cref{cor:euclid_lemma}, and $h$ has at most $p - 1 - d$ roots by \cref{thm:lagrange_roots}. So $x^d - 1$ has at least $d$ roots, and at most $d$ by \cref{thm:lagrange_roots} again.
\end{proof}

\begin{theorem}\label{thm:primitive_root_prime}
    Let $p$ be prime. For each $d \mid p - 1$ there are exactly $\phi(d)$ elements of order $d$ modulo $p$. In particular, there are $\phi(p-1) \ge 1$ primitive roots modulo $p$.
\end{theorem}
\begin{proof}
    For $d \mid p - 1$, let $\psi(d)$ be the number of $a \in \{1, \dots, p-1\}$ with $\ord_p(a) = d$. By \cref{prop:order_divides}, the solutions of $x^d \equiv 1$ are the elements whose order divides $d$, so \cref{lem:roots_xd_minus_1} gives
    \begin{equation*}
        \sum_{e \mid d} \psi(e) = d \qquad \text{for every } d \mid p - 1.
    \end{equation*}
    We show $\psi(d) = \phi(d)$ by strong induction on $d$. For $d = 1$ both are $1$. If the claim holds for all proper divisors of $d$, then \cref{prop:sum_phi_divisors} gives
    \begin{equation*}
        \psi(d) = d - \sum_{\substack{e \mid d \\ e < d}} \psi(e) = d - \sum_{\substack{e \mid d \\ e < d}} \phi(e) = \phi(d).
    \end{equation*}
    Taking $d = p - 1$ gives the number of primitive roots.
\end{proof}
For example, $2$ is a primitive root modulo $11$: its powers modulo $11$ are $2, 4, 8, 5, 10, 9, 7, 3, 6, 1$.

\section{Primitive Roots Modulo Prime Powers}

\begin{lemma}\label{lem:primitive_root_lift_choice}
    Let $p$ be an odd prime and $g$ a primitive root modulo $p$. Then $g$ or $g + p$ is a primitive root modulo $p$ that satisfies $x^{p-1} \not\equiv 1 \pmod{p^2}$.
\end{lemma}
\begin{proof}
    Both $g$ and $g + p$ are primitive roots modulo $p$. Suppose $g^{p-1} \equiv 1 \pmod{p^2}$. By the binomial theorem,
    \begin{equation*}
        (g + p)^{p-1} \equiv g^{p-1} + (p-1)p\,g^{p-2} \equiv 1 - p\,g^{p-2} \pmod{p^2},
    \end{equation*}
    and $p\,g^{p-2} \not\equiv 0 \pmod{p^2}$ because $p \nmid g$.
\end{proof}

\begin{lemma}\label{lem:lifting_powers}
    Let $p$ be an odd prime and suppose $g^{p-1} = 1 + pt$ with $p \nmid t$. Then for every $k \ge 1$,
    \begin{equation*}
        g^{p^{k-1}(p-1)} = 1 + p^k t_k \qquad \text{with } p \nmid t_k.
    \end{equation*}
\end{lemma}
\begin{proof}
    For $k = 1$ take $t_1 = t$. If the claim holds for $k$, then
    \begin{equation*}
        g^{p^{k}(p-1)} = (1 + p^k t_k)^p = 1 + p^{k+1} t_k + \sum_{j=2}^{p} \binom{p}{j} p^{jk} t_k^j.
    \end{equation*}
    For $2 \le j \le p - 1$, $p \mid \binom{p}{j}$, so the term is divisible by $p^{2k+1}$. The last term is divisible by $p^{pk}$, and $p \ge 3$. Since $2k + 1 \ge k + 2$ and $pk \ge k + 2$ for $k \ge 1$, the whole sum is divisible by $p^{k+2}$. Hence $g^{p^k(p-1)} = 1 + p^{k+1}t_{k+1}$ with $t_{k+1} \equiv t_k \not\equiv 0 \pmod p$.
\end{proof}

\begin{theorem}\label{thm:primitive_root_prime_power}
    Let $p$ be an odd prime and $k \ge 1$. There exists a primitive root modulo $p^k$. If $g$ is a primitive root modulo $p^k$, then whichever of $g$ and $g + p^k$ is odd is a primitive root modulo $2p^k$.
\end{theorem}
\begin{proof}
    Choose $g$ as in \cref{lem:primitive_root_lift_choice}, so that \cref{lem:lifting_powers} applies, and let $n = \ord_{p^k}(g)$. Then $n \mid \phi(p^k) = p^{k-1}(p-1)$. Since $g^n \equiv 1 \pmod p$ and $g$ is a primitive root modulo $p$, we have $p - 1 \mid n$. Hence $n = p^j(p-1)$ for some $0 \le j \le k - 1$.

    If $j \le k - 2$, then $n \mid p^{k-2}(p-1)$ and $g^{p^{k-2}(p-1)} \equiv 1 \pmod{p^k}$. This contradicts \cref{lem:lifting_powers}, which gives $g^{p^{k-2}(p-1)} = 1 + p^{k-1}t_{k-1}$ with $p \nmid t_{k-1}$. So $j = k - 1$ and $g$ is a primitive root modulo $p^k$.

    For the second claim, let $h \in \{g, g + p^k\}$ be odd. Then $h \equiv g \pmod{p^k}$ and $\phi(2p^k) = \phi(p^k)$ by \cref{cor:phi_multiplicative}. Since $h$ is odd, $h^n - 1$ is even for every $n$, so by \cref{thm:crt} $h^n \equiv 1 \pmod{2p^k}$ if and only if $h^n \equiv 1 \pmod{p^k}$. Hence $\ord_{2p^k}(h) = \ord_{p^k}(g) = \phi(2p^k)$.
\end{proof}

\begin{proposition}\label{prop:no_primitive_root_power_of_2}
    For $k \ge 3$ and every odd $a$, $a^{2^{k-2}} \equiv 1 \pmod{2^k}$. In particular, there is no primitive root modulo $2^k$ for $k \ge 3$, while $1$ and $3$ are primitive roots modulo $2$ and $4$.
\end{proposition}
\begin{proof}
    For $k = 3$: writing $a = 2b + 1$, we get $a^2 = 4b(b+1) + 1 \equiv 1 \pmod 8$, since $b(b+1)$ is even. If $a^{2^{k-2}} = 1 + 2^k s$, then squaring gives $a^{2^{k-1}} = 1 + 2^{k+1}s + 2^{2k}s^2 \equiv 1 \pmod{2^{k+1}}$, which is the claim for $k + 1$. Since $2^{k-2} < \phi(2^k) = 2^{k-1}$, no odd $a$ has order $\phi(2^k)$.
\end{proof}

\begin{theorem}\label{thm:primitive_root_classification}
    A primitive root modulo $m \ge 1$ exists if and only if $m \in \{1, 2, 4\}$, $m = p^k$ or $m = 2p^k$ with $p$ an odd prime and $k \ge 1$.
\end{theorem}
\begin{proof}
    For $m = 1, 2, 4$ take $g = 1, 1, 3$. The cases $p^k$ and $2p^k$ are \cref{thm:primitive_root_prime_power}.

    Conversely, first note that $\phi(n)$ is even for $n > 2$: if an odd prime $p$ divides $n$ then $p - 1 \mid \phi(n)$, and otherwise $n = 2^k$ with $k \ge 2$ and $\phi(n) = 2^{k-1}$. Suppose $m = m_1m_2$ with $(m_1, m_2) = 1$ and $m_1, m_2 > 2$. Let $L$ be the least common multiple of $\phi(m_1)$ and $\phi(m_2)$. Since both are even, $(\phi(m_1),\phi(m_2)) \ge 2$, and \cref{cor:divisors_valuations} gives
    \begin{equation*}
        L = \frac{\phi(m_1)\phi(m_2)}{(\phi(m_1),\phi(m_2))} \le \frac{\phi(m)}{2}.
    \end{equation*}
    For $(a,m) = 1$, \cref{thm:euler_theorem} gives $a^L \equiv 1$ modulo $m_1$ and modulo $m_2$, hence modulo $m$. So $\ord_m(a) \le \phi(m)/2$ and there is no primitive root.

    If $m$ is not in the list, then either $m = 2^k$ with $k \ge 3$, which is \cref{prop:no_primitive_root_power_of_2}, or $m$ can be split as above: take $m_1 = p^{v_p(m)}$ for an odd prime $p \mid m$ and $m_2 = m/m_1$. Then $m_1 > 2$, and $m_2 > 2$ because $m_2$ is divisible by another odd prime or by $4$.
\end{proof}

